% Figura: problema 4.4. Perfiles axiales de E (unid. P/eps0), D y P (unid. P) para
% tres relaciones de aspecto a/l.  Funciones analíticas: \curvaEeje, \curvaDeje (preámbulo).
\begin{figure}[t]
\centering
\newcommand{\icono}[2]{\draw[texto, fill=dielecB, line width=0.35pt]
   (rel axis cs:0.86,0.82) ++(-#1/2,-#2/2) rectangle ++(#1,#2);}
\begin{tikzpicture}
\begin{groupplot}[group style={group size=3 by 2, horizontal sep=0.95cm, vertical sep=1.05cm},
   perfil, width=5.55cm, height=4.0cm]
  % --- fila superior: E y P
  \nextgroupplot[ymin=-1.15, ymax=1.15, ytick={-1,-0.5,0,0.5,1}, ylabel={$E,\ P$}]
    \curvaPeje \curvaEeje{1}
    \node[etiqf, font=\tiny, anchor=north west] at (rel axis cs:0.02,0.98) {$a/\ell=1$};
    \icono{0.30cm}{0.30cm}
  \nextgroupplot[ymin=-1.15, ymax=1.15, ytick={-1,-0.5,0,0.5,1}]
    \curvaPeje \curvaEeje{50}
    \node[etiqf, font=\tiny, anchor=north west] at (rel axis cs:0.02,0.98) {$a/\ell=50$};
    \icono{0.05cm}{0.5cm}
  \nextgroupplot[ymin=-1.15, ymax=1.15, ytick={-1,-0.5,0,0.5,1}]
    \curvaPeje \curvaEeje{0.1}
    \node[etiqf, font=\tiny, anchor=north west] at (rel axis cs:0.02,0.98) {$a/\ell=0{,}1$};
    \icono{0.55cm}{0.06cm}
  % --- fila inferior: D y P
  \nextgroupplot[ymin=-0.08, ymax=1.15, ytick={0,0.5,1}, ylabel={$D,\ P$}]
    \curvaPeje \curvaDeje{1}
  \nextgroupplot[ymin=-0.08, ymax=1.15, ytick={0,0.5,1}]
    \curvaPeje \curvaDeje{50}
  \nextgroupplot[ymin=-0.08, ymax=1.15, ytick={0,0.5,1}]
    \curvaPeje \curvaDeje{0.1}
\end{groupplot}
\end{tikzpicture}
\caption{\Cref{prob:cilindro-polarizado}. Campos sobre el eje del electrete cilíndrico: $E_z$ en unidades de $P/\eo$ (continua, cian),
$D_z$ en unidades de $P$ (continua, violeta) y $P$ (trazos). Arriba: $E_z$ salta $P/\eo$ en cada cara y dentro se
opone a $\vP$. Abajo: $D_z$ es continuo (no hay carga libre) y coincide con la polarización solo en el límite del
cilindro largo. Columnas: cilindro equilátero, disco delgado ($E\simeq-P/\eo$, $D\simeq0$) y cilindro largo
($E\simeq0$, $D\simeq P$ lejos de los extremos).}
\label{fig:p44-perfiles}
\end{figure}
